\chapter{Experimental Setup}
\label{chap:experimental-setup}

Chapter 3 defined QCLand's architecture, training objective, and
inference procedure. This chapter specifies how the framework is
evaluated. It describes the datasets and their released partitions,
defines the evaluation metric, identifies the compared models,
establishes the shared protocol for the final four-model comparison,
and separates the BPD development experiments from the cross-model
evaluation.


% =================================================================
% 4.1 Datasets and Data Splits
% =================================================================

\section{Datasets and Data Splits}
\label{sec:setup-datasets}

All experiments use publicly released, de-identified fetal ultrasound
data from the Multicentre Fetal Biometry benchmark
\parencite{divece2025multicentre}. No new patient data or annotations
were collected. The pooled release was used under its CC~BY-NC-SA~4.0
licence for non-commercial academic research; original images are not
redistributed with this thesis. Ethical approvals for the constituent
datasets are recorded in \Cref{app:ethics}.


\subsection{UCL Single-Site Dataset}
\label{sec:setup-ucl}

Experiments cover the five fetal measurements introduced in
Chapter 1, with BPD/OFD sharing the Head images, TAD/APAD sharing
the Abdomen images, and FL using the Femur images. Each measurement
is defined by two anatomical landmarks per image. The experiments
retain the official Train/Test partitions released by
\textcite{divece2025multicentre} without modification.

For early stopping and checkpoint selection, a portion of each
official training pool is held out for internal validation. The split
is made over filename-prefix groups, with all frames sharing a leading
numeric prefix assigned to the same partition, using a fixed split
seed of 42. The last
$\lceil 0.1 \times n_{\text{groups}}\rceil$ groups form the
validation set. A filename-prefix group therefore cannot occur in
both the optimisation and validation sets. The exact image counts are
reported in \Cref{tab:full-dataset-sizes}.


\subsection{Pooled Multicentre Dataset}
\label{sec:setup-multicentre}

The Multicentre release \parencite{divece2025multicentre} pools UCL
with two further sources: FP, collected at two hospitals in Barcelona,
and HC18, collected at a clinical site in the Netherlands
\parencite{vandenheuvel2018hc18}. Together, the release spans four
clinical sites and seven ultrasound device models.
\Cref{tab:multicentre-per-source} shows the per-source composition.

\begin{table}[htbp]
  \centering
  \caption{Per-source image counts (Train / Test) in the pooled
    Multicentre release, by anatomy.}
  \label{tab:multicentre-per-source}
  \begin{tabular}{lccc}
    \toprule
    Anatomy & FP & HC18 & UCL \\
    \midrule
    Head     & 757 / 880 & 737 / 262 & 110 / 49 \\
    Abdomen  & 568 / 125 & --        & 94 / 36  \\
    Femur    & 437 / 323 & --        & 96 / 39  \\
    \bottomrule
  \end{tabular}
\end{table}

Internal validation groups are sampled from each pooled training
split using a fixed seed of 42 and source-aware filename grouping.
The released Train/Test partition itself is retained unchanged.
Because the pooled benchmark concatenates source-specific partitions
rather than holding out an entire clinical site, the Multicentre
experiments assess performance within the released multi-site
benchmark rather than unseen-site transfer.

\begin{table}[htbp]
  \centering
  \caption{Exact image counts for every dataset configuration used
    in this thesis.}
  \label{tab:full-dataset-sizes}
  \begin{tabular}{lcccc}
    \toprule
    Dataset (task) & Official pool & Optimisation & Validation & Test \\
    \midrule
    UCL Head (\BPD/\OFD)     & 110  & 100  & 10  & 49 \\
    UCL Abdomen (\TAD/\APAD) & 94   & 85   & 9   & 36 \\
    UCL Femur (\FL)          & 96   & 84   & 12  & 39 \\
    Multicentre Head         & 1604 & 1453 & 151 & 1191 \\
    Multicentre Abdomen      & 662  & 590  & 72  & 161 \\
    Multicentre Femur        & 533  & 481  & 52  & 362 \\
    \bottomrule
  \end{tabular}
\end{table}


% =================================================================
% 4.2 Evaluation Metric
% =================================================================

\section{Evaluation Metric}
\label{sec:setup-nme}

The two landmarks defining each fetal measurement form an unordered
pair: exchanging their labels does not change the underlying clinical
measurement. The final comparison therefore uses
permutation-invariant normalised mean error (PI-NME), allowing all
methods to be evaluated independently of their internal landmark
ordering conventions.

For predicted landmarks $\hat{p}_0,\hat{p}_1$ and ground-truth
landmarks $p_0,p_1$, the two possible correspondences are

\begin{align}
  d_{\text{direct}}
  &=
  \lVert \hat{p}_0 - p_0 \rVert_2
  +
  \lVert \hat{p}_1 - p_1 \rVert_2,
  \\
  d_{\text{crossed}}
  &=
  \lVert \hat{p}_0 - p_1 \rVert_2
  +
  \lVert \hat{p}_1 - p_0 \rVert_2.
\end{align}

With

\begin{equation}
  D = \lVert p_0 - p_1 \rVert_2,
\end{equation}

PI-NME is defined as

\begin{equation}
  \label{eq:pi-nme}
  \mathrm{PI\text{-}NME}
  =
  \frac{
    \min(d_{\text{direct}},\,d_{\text{crossed}})
  }{2D}.
\end{equation}

All distances are computed in the original image coordinate system.
Predicted coordinates are first mapped back from model-input space
using each method's own inverse geometric transform, and only then
are the two possible correspondences and final PI-NME evaluated.
This is important because the original ultrasound frames are
generally non-square, whereas the compared models operate on square
resized inputs; computing Euclidean errors directly in resized
coordinates would distort the relative contributions of the
horizontal and vertical axes.

This external PI-NME definition is applied identically to QCLand,
HRNet-W18, and RTMPose-s in the final comparison. HRNet-W18 retains
its DOD-based landmark-ordering mechanism, while the other methods
retain their own internal assignment conventions; these internal
choices do not alter the common permutation-invariant metric used
for final evaluation.


% =================================================================
% 4.3 Compared Methods and Baselines
% =================================================================

\section{Compared Methods and Baselines}
\label{sec:setup-baselines}

The final comparison includes four models: QCLand-DINOv2,
QCLand-DINOv3, HRNet-W18, and RTMPose-s. The two QCLand variants
share the architecture defined in Chapter 3 and differ only in the
pretrained backbone.


\subsection{Locally Reproduced HRNet-W18}

HRNet-W18 is reproduced from the original authors' implementation
using the released training protocol. The upstream procedure monitors
the official test partition during training rather than constructing
a separate internal validation set, and the final-state checkpoint is
used for evaluation. This differs from QCLand and RTMPose-s, which
reserve part of the training pool for internal validation. Full
reproduction details and software compatibility modifications are
documented in \Cref{app:reproducibility}.

Published BiometryNet results from
\textcite{divece2025multicentre} are retained only as contextual
references rather than protocol-matched comparisons, because their
complete training and evaluation conditions cannot be independently
matched to the present experiments.


\subsection{RTMPose-s}
\label{sec:setup-rtmpose}

RTMPose-s \parencite{jiang2023rtmpose} is adapted as a two-landmark
model trained independently for each dataset--measurement setting,
matching QCLand's and HRNet-W18's per-measurement organisation. Its
native CSPNeXt-s backbone and RTMCC/SimCC formulation are retained;
input resolution, data partitions, seeds, and external evaluation
are aligned with the final comparison protocol.

RTMPose-s trains for a fixed 200 epochs with no early stopping.
Its internal validation measurements are diagnostic only, and the
final training state is reported. The full adaptation protocol,
including codec configuration and software versions, is documented
in \Cref{app:reproducibility}.


% =================================================================
% 4.4 Final Four-Model Comparison Protocol
% =================================================================

\section{Final Four-Model Comparison Protocol}
\label{sec:setup-final-comparison}
\label{sec:setup-four-model}

The four models are compared across all ten dataset--measurement
settings: five measurements on UCL and the same five measurements
on Multicentre. \Cref{tab:comparison-conditions} summarises the
shared and method-specific evaluation conditions.

\begin{table}[htbp]
  \centering
  \caption[Four-model comparison protocol.]
  {Harmonised and method-specific conditions for the final
    four-model comparison.}
  \label{tab:comparison-conditions}
  \small
  \begin{tabularx}{\textwidth}{
    @{}>{\raggedright\arraybackslash}p{0.22\textwidth}
    *{4}{>{\centering\arraybackslash}X}@{}}
    \toprule
    Condition &
    QCLand-v2 &
    QCLand-v3 &
    HRNet-W18 &
    RTMPose-s \\
    \midrule
    Test partition
      & released & released & released & released \\
    Input resolution
      & $512^2$ & $512^2$ & $512^2$ & $512^2$ \\
    Training seeds
      & 5 named & 5 named & 5 named & 5 named \\
    External metric
      & PI-NME & PI-NME & PI-NME & PI-NME \\
    Coordinate space
      & original & original & original & original \\
    \midrule
    Training termination
      & validation-based
      & validation-based
      & upstream protocol
      & fixed 200 epochs \\
    Validation use
      & stopping only
      & stopping only
      & test monitored
      & diagnostic internal \\
    Reported checkpoint
      & terminal state
      & terminal state
      & final state
      & final state \\
    \bottomrule
  \end{tabularx}
\end{table}

Training-time validation and termination rules differ across methods
and are not harmonised. QCLand uses an internal validation split to
determine when training terminates. Validation determines the stopping
epoch but not which checkpoint is reported: the reported model is the
terminal checkpoint saved when training stops, rather than the
checkpoint with the lowest validation PI-NME.

HRNet-W18 retains the upstream training protocol, in which the
official test partition is monitored during training while the
upstream final-state checkpoint is used for evaluation. RTMPose-s
trains for a fixed 200 epochs; its internal validation measurements
are diagnostic only, and the final training state is reported.

The comparison therefore harmonises the reported terminal or final
model state rather than the mechanism by which each training run
reaches that state.


\subsection{Common Test-Image Subsets}
\label{sec:results-common-subset}

Not every method's data pipeline retains identical images for a given
dataset--measurement setting. HRNet-W18's upstream loader excludes
rows with negative or missing landmark coordinates. Every formal
comparison therefore uses the filename intersection shared across all
four methods and all five training seeds. The exact retained counts
per setting are reported alongside the results in Chapter 5.


\subsection{Five-Seed Training}
\label{sec:setup-multiseed}

All four methods use five fixed training seeds: 42, 0, 123, 2024,
and 3407. The seed controls weight initialisation, data order, and
stochastic augmentation but does not change the train/validation
partition.

Results are reported as the mean and sample standard deviation
(ddof\,=\,1) across the five independently trained models. The
reported $\pm$ therefore represents run-to-run variability rather
than a confidence interval.


\subsection{Paired Statistical Analysis}

For each of the ten dataset--measurement settings and six method pairs
per setting, a per-image score is first formed by averaging PI-NME
across the five training seeds on the common image subset. Pairwise
differences between methods are then formed from these per-image
five-seed means and analysed using two complementary procedures:

\begin{itemize}

  \item a paired bootstrap with 20{,}000 resamples, reporting the
    2.5th and 97.5th percentiles of the resampled mean difference;

  \item a two-sided Wilcoxon signed-rank test using Pratt zero
    handling, with Holm correction applied both within each
    setting's six pairwise contrasts and globally across all sixty
    contrasts (Holm-60).

\end{itemize}

Negative pairwise differences favour method A when the contrast is
defined as method A minus method B. The bootstrap confidence interval
and Wilcoxon test characterise complementary aspects of the paired
difference distribution and are both reported when their conclusions
differ.


% =================================================================
% 4.5 BPD Development Experiments
% =================================================================

\section{BPD Development Experiments}
\label{sec:setup-bpd-development}


\subsection{Scope of the BPD Development Study}
\label{sec:setup-evidence-scope}

BPD is the sole method-development task. The development sequence
comprises target-width screening, loss screening, the
decoder/FPN/UDP component chain, augmentation analysis, and
input-resolution sensitivity. No other measurement repeats this
complete sequence. Results on the remaining four measurements
therefore test the adopted pipeline as a whole rather than isolating
the independent contribution of every component.

The BPD development study and the final cross-model comparison
therefore provide distinct forms of evidence: the former examines
architectural and training choices during model development, whereas
the latter evaluates the adopted framework across measurements,
backbones, and dataset settings.


\subsection{Matched Controls}

The target-$\sigma$ screen is a single-seed exploratory comparison
using seed 2024. The loss screen and the main component chain use the
five predefined seeds and PI-NME from the reported terminal
checkpoints.

The augmentation follow-up uses matched five-seed paired analysis on
the 49 common UCL BPD test images, with 20{,}000 bootstrap resamples
and a two-sided Wilcoxon signed-rank test. The input-resolution
analysis uses a separate matched five-seed configuration and does not
supply the QCLand-DINOv2 entry in the final four-model comparison.


% =================================================================
% 4.6 QCLand Experimental Configuration
% =================================================================

\section{QCLand Experimental Configuration}
\label{sec:setup-default-config}

\Cref{tab:default-config} summarises the adopted hyperparameters and
training settings for QCLand. Architectural components
(DeconvHeadV2, multi-depth fusion, the hybrid loss, and coordinate
handling) are defined in Chapter 3; this table specifies their
numerical settings and the training recipe.

\begin{table}[htbp]
  \centering
  \caption[QCLand experimental configuration.]
  {QCLand experimental configuration for fetal landmark localisation.
    Rotation-and-scale augmentation is enabled for the final
    five-measurement UCL evaluation and all Multicentre experiments;
    the unaugmented configuration is used only in the BPD development
    chain where explicitly stated.}
  \label{tab:default-config}
  \small
  \begin{tabular}{@{}p{0.32\textwidth}p{0.62\textwidth}@{}}
    \toprule
    Setting & Value \\
    \midrule
    Backbone
      & DINOv2 ViT-S/14 (patch kernel resampled to
        $16{\times}16$) or DINOv3 ViT-S/16 (native) \\
    Input resolution
      & $512 \times 512$ (non-aspect-preserving) \\
    Heatmap resolution
      & $64 \times 64$ \\
    Gaussian $\sigma$
      & 4.0 \\
    Loss
      & Weighted MSE ($\alpha{=}5.0$) + soft-argmax coordinate
        L1 ($\tau{=}10$, $\lambda_{\text{coord}}{=}0.1$) \\
    Optimiser
      & AdamW; base LR $10^{-4}$; weight decay 0.05 \\
    Layer-wise LR decay
      & $\rho = 0.8$ \\
    Warmup
      & Head parameters: 15 steps; backbone held at zero for
        15 steps, followed by 30-step warmup \\
    LR schedule
      & Polynomial decay (power 0.9) after warmup \\
    Batch size
      & 16 \\
    Max epochs
      & 200 \\
    Validation interval
      & Every 5 epochs \\
    Validation-based termination
      & 20 consecutive validation checks without
        $>0.005$ NME improvement \\
    Base augmentation
      & Horizontal flip and brightness/contrast/saturation jitter,
        each with $p{=}0.5$ \\
    Final augmentation
      & Rotation $\pm30^\circ$ ($p{=}0.6$) and scale
        $\pm25\%$; skipped if a landmark would leave the image \\
    \bottomrule
  \end{tabular}
\end{table}


% =================================================================
% 4.7 Implementation and Reproducibility
% =================================================================

\section{Implementation and Reproducibility}
\label{sec:setup-implementation}

Full details of the software and hardware environments, seeding
methodology, baseline reproduction procedures, statistical
implementation, and per-run configuration provenance are provided in
\Cref{app:reproducibility}. The implementation and reproduction
scripts will be made publicly available at
\url{https://github.com/Flora020703/eomt-landmark-detection}
upon submission.

With the datasets, evaluation metric, comparison protocol, and
development-study scope defined, Chapter 5 reports the corresponding
experimental results.
